\icode{E-MD-MU-L3-261115-004}
图甲所示为电梯的结构示意图，配重和轿厢质量均为 \(M_0 =\) \,\qty{500}{kg}，\(A\) 与电机（未画出）相连，缆绳质量以及其与滑轮 \(B\) 间的摩擦均可忽略。家住 21 楼的小明要爸爸送他上幼儿园，电梯下行过程经加速、匀速、减速三个阶段，最后停在一楼。已知小明质量为 \(m_1 =\) \,\qty{20}{kg}，爸爸的质量为 \(m_2 =\) \,\qty{60}{kg}，小明对轿厢底的压力随时间变化的关系图像如图乙所示，重力加速度 \(g =\) \,\qty{10}{m/s^2}。则下列说法正确的是\paren[]
	
\begin{tikzpicture}[font=\scriptsize]
	\draw (.1,3)node[above left]{\(A\)} circle(.1) (1.9,3)node[above right]{\(B\)} circle(.1);
	\fill (.1,3) circle(.05) (1.9,3) circle(.05);
	\draw[magenta] (0,.7)--(0,3) (.1,3.1)--(1.9,3.1) (2,3)--(2,2);
	\draw[fill=orange!50] (-.25,0) rectangle(.25,.7);
	\draw[fill=gray!50] (1.85,2) rectangle(2.15,1.5);
	\node at(.25,.35)[right]{轿厢};
	\node at(1.85,1.75)[left]{配重};
	\node at(1,-.6){图甲};
\end{tikzpicture}
\hspace{3em}
\begin{tikzpicture}[>=stealth,font=\scriptsize]
	\draw[->,blue] (-.4,0)--(5,0)node[below]{\(t\)/\qty{}{s}};
	\draw[->,blue] (0,-.4)--(0,.2)--(-.2,.25)--(.2,.3)--(0,.35)--(0,3)node[left]{\(F\)/\qty{}{N}};
	\draw[magenta] (0,1)--(.4,1) (0.4,1.5)--(3.2,1.5) (3.2,2)--(3.6,2) (3.6,1.5)--(4.5,1.5);
	\draw[dashed,cyan] (0,1.5)--(.4,1.5)--(.4,0) (0,2)--(3.2,2)--(3.2,0) (3.6,2)--(3.6,0);
	\foreach \y/\n in {1/190,1.5/200,2/210} 
		\node[blue] at(0,\y)[left]{\(\n\)};
	\foreach \x/\n in {.4/4,3.2/32,3.6/36}
		\node[blue] at(\x,0)[below]{\(\n\)};
	\node[blue] at(-.18,-.18){\(0\)};
	\node at(2.7,-.6){图乙};
\end{tikzpicture}
\begin{choices}
	\item 小明家所在楼栋楼层的平均高度为\qty{3}{m}
	\item 加速阶段，缆绳对配重的拉力大小为\qty{500}{N}
	\item 整个过程中，电动机的最大功率为\qty{2680}{W}
	\item 整个过程中，电动机对外做的功大小为\qty{51200}{J}
\end{choices}
\begin{answer}
    ACD\qquad 提示：
\end{answer}